% Generated by build2/mktab.py from the R pipeline output (tab3_components.tex); layout only.
\begin{table}[!htbp]
\centering
\begin{threeparttable}
\caption{Retained principal components (main analysis)}
\label{tab:pcs}
\footnotesize\setlength{\tabcolsep}{4pt}
\begin{tabular}{@{}l S[table-format=2.1] S[table-format=2.1] S[table-format=2.1] S[table-format=2.0] *{5}{S[table-format=-1.2]}@{}}
\toprule
 & {Variance} & {Cumulative} & {Landing} & {Peak} & \multicolumn{5}{c}{Pearson $r$ with} \\
\cmidrule(l){6-10}
Component & {(\%)} & {(\%)} & {share (\%)} & {at (\%)} & {LPF} & {JH} & {CT} & {CMD} & {LD} \\
\midrule
PC1 & 37.3 & 37.3 & 99.7 & 40 & 0.39 & 0.13 & -0.04 & 0.01 & -0.75 \\
PC2 & 16.0 & 53.3 & 99.9 & 26 & 0.22 & 0.13 & 0.02 & 0.00 & -0.14 \\
PC3 & 12.6 & 65.9 & 99.5 & 16 & 0.39 & 0.34 & 0.00 & -0.14 & -0.32 \\
PC4 & 8.9 & 74.8 & 99.3 & 9 & 0.45 & 0.24 & 0.15 & -0.18 & -0.04 \\
\bottomrule
\end{tabular}
\begin{tablenotes}[flushleft]\footnotesize
\item \textit{Note.} Landing share = percentage of the squared loading vector that lies in the landing block (the rest is in take-off). Peak at = position of the largest absolute loading, in \% of the landing phase. Signs are set so that the largest absolute loading is positive. LPF = peak landing force relative to body weight; JH = jump height; CT = contraction time; CMD = countermovement depth; LD = landing duration (touch-down to lowest COM).
\end{tablenotes}
\end{threeparttable}
\end{table}
